\chapter{Familial Hypercholesterolemia}
\index{Familial Hypercholesterolemia}
\label{sec:Familial Hypercholesterolemia}

\section{Overview}
Familial hypercholesterolemia (FH) is an autosomal codominant disorder of LDL metabolism caused by loss-of-function mutations in the LDLR gene (LDL receptor, 90\% of cases), APOB (apolipoprotein B-100), or gain-of-function mutations in PCSK9. It is the most common monogenic lipid disorder, with a prevalence of 1 in 250 for heterozygous FH (HeFH) and 1 in 300,000 for homozygous FH (HoFH).
\section{Causes}
This condition happens because of impaired liver clearance of LDL particles.     
\section{Risk Factors}

\begin{itemize}[noitemsep]
  \item Genetic predisposition and family history
  \item Advancing age
  \item Lifestyle factors (diet, exercise, smoking, alcohol)
  \item Environmental exposures
  \item Underlying medical conditions
  \item Medication use
\end{itemize}
\section{Signs and Symptoms}

\subsection{Common symptoms}
\begin{itemize}[noitemsep]
  \item You may experience tendon xanthomas (Achilles, extensor tendons of the hands), tuberous xanthomas, xanthelasmas, corneal arcus before age 45, and premature atherosclerotic cardiovascular disease
  \item Homozygous FH presents in childhood with LDL-C greater than 500 mg/dL, widespread xanthomas, and coronary atherosclerosis before age 20.
  \item Pain or discomfort in the affected area
  \end{itemize}

\subsection{Emergency warning signs}
\begin{itemize}[noitemsep]
  \item Severe or worsening symptoms of familial hypercholesterolemia require immediate medical evaluation
\end{itemize}
\section{Progression and Stages}
The condition may progress through different stages of severity over time. Early stages often involve mild or subtle symptoms that may be overlooked. Moderate stages involve more noticeable symptoms affecting daily function. Advanced stages may involve significant impairment and complications.
\section{Complications}
\begin{itemize}[noitemsep]
  \item Disease progression leading to worsening symptoms
  \item Secondary infections or injuries
  \item Side effects of treatment
  \item Reduced quality of life
  \item Emotional and psychological impact
\end{itemize}
\section{Diagnosis}
Doctors diagnose this using clinical criteria (Dutch Lipid Clinic Network or Simon Broome criteria) supported by genetic testing. Comprehensive medical history and physical examination Laboratory tests to assess organ function and rule out other conditions Imaging studies to visualize affected structures Specialized testing depending on the affected system Consultation with relevant specialists Monitoring of disease progression over time

\begin{itemize}[noitemsep]
  \item Comprehensive medical history and physical examination
  \item Laboratory tests to assess organ function
  \item Imaging studies to visualize affected structures
  \item Specialized testing depending on the affected system
  \item Consultation with relevant specialists
\end{itemize}
\section{Prevention}
\begin{itemize}[noitemsep]
  \item Maintain a healthy lifestyle with balanced diet and exercise
  \item Avoid known risk factors
  \item Attend regular medical check-ups for early detection
  \item Follow recommended screening guidelines
\end{itemize}
\section{Treatment Options}
Treatment options include high-dose statins with ezetimibe as first-line for HeFH

\begin{itemize}[noitemsep]
  \item PCSK9 inhibitors for inadequate response
  \item For HoFH, all available agents and LDL apheresis are needed.
  \item Medications to manage symptoms and address underlying mechanisms
  \item Lifestyle modifications and self-care measures
  \item Physical therapy and rehabilitation
  \item Surgical intervention when indicated
  \item Regular monitoring and follow-up care
\end{itemize}
\section{Living with It}
\begin{itemize}[noitemsep]
  \item Follow your treatment plan as prescribed
  \item Maintain regular follow-up with your healthcare provider
  \item Make necessary lifestyle adjustments
  \item Seek support from family, friends, or support groups
  \item Stay informed about your condition
\end{itemize}
\section{Outlook}
Most people recover well with early diagnosis and treatment. The outlook varies depending on the specific condition, severity at diagnosis, and response to treatment. Early intervention generally leads to better outcomes. Many conditions can be effectively managed with appropriate treatment. Regular follow-up care is important for monitoring and treatment adjustment.
